% ==========================================
% SEZIONI DEL VERBALE
% ==========================================
\newpage
\vspace{0.8cm}

\section{Ordine del giorno}

\subsection{Presentazione e approccio al progetto}
Durante la fase di apertura dell'incontro, il team di progetto ha effettuato le presentazioni generali con il referente aziendale. Quest'ultimo ha esplicitato l'approccio dell'azienda nei confronti dello sviluppo: da un lato viene garantita un'ampia flessibilità lato requisiti e scelte implementative minime, dall'altro si richiede la massima formalità e il rigoroso rispetto dei vincoli architetturali, documentali e metodologici previsti dal corso di studi universitario.

\subsection{Limiti del piano Hobby di Vercel}
Per la fase iniziale di prototipazione, il team utilizzerà l'ambiente di deployment fornito dal piano \textit{Vercel Hobby}. È stato stabilito che, qualora venissero rilevate limitazioni operative (es. restrizioni sui tempi di esecuzione, sulle risorse computazionali o sul numero di esecuzioni schedulate) oppure rallentamenti prestazionali, l'azienda provvederà direttamente a fornire le credenziali per un account aziendale sul piano \textit{Vercel Pro}.

\subsection{Gestione token Shopify e sicurezza delle credenziali}
L'analisi si è concentrata sulle nuove direttive applicative di Shopify (in vigore dal 2026), che vietano la creazione di app legacy a favore di un meccanismo di autenticazione dinamico gestito tramite session token:
\begin{itemize}
    \item \textbf{Token ed Errore 401:} I token di sessione generati durano 24 ore. L'applicazione deve intercettare l'errore HTTP 401 per rigenerare il token automaticamente.
    \item \textbf{Caching e Concorrenza:} Il token deve essere salvato in una cache a basso tempo di latenza (es. Vercel Runtime Cache KV o tabella a database) per evitare la rigenerazione simultanea e concorrente di più token ad ogni scansione.
    \item \textbf{Sicurezza Client Secret:} Il Client Secret non deve mai essere salvato in chiaro né esposto al frontend o nelle risposte API. In futuro si prevede di gestire anche la rotazione periodica del secret.
\end{itemize}

\subsection{Architettura multi-tenancy, Login con Postgres}
La definizione del modello operativo multi-tenant costituisce uno dei requisiti architetturali primari:
\begin{itemize}
    \item \textbf{Isolamento Tenant:} L'applicazione deve supportare più organizzazioni senza sovrapposizione di dati. Ogni utente appartiene a un'unica organizzazione.
    \item \textbf{Row Level Security (RLS):} È stato concordato di delegare l'isolamento al database tramite RLS su PostgreSQL/Supabase. La vista dell'applicazione effettuerà la query inserendo il contesto dell'utente loggato, garantendo che le righe restituite siano solo quelle di pertinenza della propria organizzazione.
    \end{itemize}
È stata concordata come prima milestone fondamentale l'implementazione dell'autenticazione di base (modulo di login, Hello World e tabelle organizzazione/utenti con RLS abilitata), da presentare e discutere nel prossimo incontro.

\subsection{Schedulazione task su Vercel}
È stato chiarito che i tradizionali \textit{Cron Jobs} sono stati ufficialmente deprecati dalla piattaforma Vercel in favore dei nuovi \textit{Vercel Schedules}. L'architettura di schedulazione concordata prevede un'interazione disaccoppiata tramite API:
\begin{itemize}
    \item La \textit{Vercel Schedule} ha il solo compito di effettuare una chiamata API ad evento/orario per attivare l'esecuzione del workflow.
    \item Il \textit{Vercel Workflow} prende in carico l'elaborazione ed esegue gli step operativi in sequenza (es. estrazione dei prodotti via API, trasformazione dati, persistenza).
\end{itemize}

\subsection{Requisiti di durabilità (RF03) e rate limiting (RF04)}
La durabilità dello stato di esecuzione è garantita nativamente dalle primitive dei Workflows di Vercel, che salvano automaticamente l'avanzamento dei vari step. Di conseguenza, non si rende necessaria l'implementazione di mock server o infrastrutture complesse per dimostrare la durabilità nel PoC.
\\ \\
Per la versione iniziale del PoC è considerata sufficiente l'adozione di un \textit{throttling} statico (introduzione di pause/delay temporizzati lato codice tra le chiamate). Nelle fasi successive, il motore analizzerà gli header HTTP di risposta restituiti dalle Shopify Admin API per calibrare dinamicamente la frequenza delle richieste ed evitare il sovraccarico delle API esterne.


\subsection{Requisiti e priorità}
In merito ai requisiti di progetto e alle relative priorità di sviluppo, sono state definite le seguenti linee guida:
\begin{itemize}
    \item \textbf{Requisiti RFD02 e RFD07:} Consigliata l'inclusione nell'MVP del requisito RFD02 (generazione descrizioni in linguaggio naturale tramite Vercel AI Gateway) e RFD07 (esportazione del grafo in JSON/XML mediante plugin) in quanto parzialmente coperti dalle funzionalità messe a disposizione da Vercel.
    \item \textbf{Falsi positivi(RF05):} Da rimandare in data successiva, si tratta di un fine-tuning da fare dopo aver raccolto dati.
    \item \textbf{Notifiche (RFO19/RFO06):} Per la prima fase è ritenuta sufficiente una vista nell'interfaccia utente che legge direttamente la tabella notifiche a DB, senza l'invio preventivo di e-mail.
    \item \textbf{Analisi del tema:} Nell'analisi del tema andrà esplicitato per ogni file processato se esso rientra nel perimetro di analisi o se viene escluso.
\end{itemize}

\subsection{Calendario prossimi incontri}
È stato definito un primo calendario di incontri che si svolgeranno a cadenza settimanale di mercoledì su Google Meet dalle ore 14 alle ore 15:
\begin{itemize}
    \item \textbf{18 Novembre 2026:} Incontro settimanale straordinario (invito a cura del team).
    \item \textbf{25 Novembre 2026:} \textit{Incontro saltato} (azienda in evento aziendale offsite).
    \item \textbf{02 Dicembre 2026 -- 09 Dicembre 2026:} Incontri ordinari di avanzamento.
    \item \textbf{15 Dicembre 2026:} Anticipo rispetto a mercoledì 16 per disponibilità referente.
    \item \textbf{07 Gennaio 2027 -- 13 Gennaio 2027:} Incontri finali pre-consegna POC (prevista per metà Gennaio).
\end{itemize}

